% =============================================================================
% INF-221 Algoritmos y Complejidad -- Tarea 1 -- Semestre 2026-2
% Autor: Kurt Wiederhold   Rol: 202473528-4
% Seccion: Introduccion (extension maxima: 1/4 de pagina)
% =============================================================================

Este trabajo estudia experimentalmente dos problemas clásicos: ordenar un
arreglo de enteros y multiplicar dos matrices cuadradas. Para el primero se
comparan \texttt{std::sort}, merge sort, quick sort y patience
sort~\cite{taocv3,clrs}; para el segundo, el método clásico de tres bucles
frente al algoritmo de Strassen~\cite{strassen1969}. Sus complejidades teóricas
son conocidas, de modo que el objetivo no es deducirlas sino contrastarlas con
el comportamiento observado sobre los casos del anexo~A del enunciado, variando
el orden inicial de la entrada, la estructura de la matriz y el dominio de los
valores. La pregunta que guía el análisis es qué esconde la notación
asintótica: dos algoritmos del mismo orden pueden diferir en un factor
constante grande, un exponente medido puede no coincidir con el teórico, y un
algoritmo asintóticamente superior puede perder en el rango de tamaños que
interesa en la práctica.
